\section{Method}

Figure~\ref{fig:pipeline} summarizes the three-stage pipeline. A small set of
normal images is encoded once to form a category-specific memory. Test images
are encoded by the same frozen model, background patches are suppressed, and
each retained patch is compared against the memory. The resulting distances
produce both a full anomaly map and an image score.

\begin{figure*}[t]
\centering
\begin{tikzpicture}[
  node distance=8mm and 10mm,
  box/.style={draw, rounded corners=2pt, align=center, minimum height=8mm,
              text width=23mm, font=\scriptsize},
  arr/.style={-{Latex[length=1.8mm]}, thick}
]
\node[box] (refs) {1/2/4/8\\normal references};
\node[box, right=of refs] (enc1) {Frozen DINOv2\\patch encoder};
\node[box, right=of enc1] (memory) {Normalized\\patch memory};
\node[box, below=of refs] (test) {Test\\image};
\node[box, right=of test] (enc2) {Frozen DINOv2\\patch encoder};
\node[box, right=of enc2] (mask) {PCA background\\mask};
\node[box, below=of memory] (knn) {Nearest-normal\\distance};
\node[box, left=of knn] (map) {Patch anomaly\\map};
\node[box, left=of map] (score) {Top-1\% score\\+ threshold};
\draw[arr] (refs) -- (enc1);
\draw[arr] (enc1) -- (memory);
\draw[arr] (test) -- (enc2);
\draw[arr] (enc2) -- (mask);
\draw[arr] (memory) -- (knn);
\draw[arr] (mask) |- (knn);
\draw[arr] (knn) -- (map);
\draw[arr] (map) -- (score);
\node[draw, dashed, fit={(enc1) (enc2)}, inner sep=3mm,
      label={[font=\scriptsize]above:Improved variant fuses blocks 9 and 11}] {};
\end{tikzpicture}
\caption{Anomaly inspection pipeline. The dashed region denotes the only
improvement over the official backbone: normalized features from two
transformer blocks are averaged before memory construction and retrieval.}
\label{fig:pipeline}
\end{figure*}

\subsection{Template-Difference Baseline}
For a category $c$, let $\mathcal{R}_{c,s,r}$ denote the $s$ normal references
selected for seed $r$. The reference set is always the same across all three
methods. Every image is resized to $448\times448$. For test image $I$ and
reference $R_j$, we build four normalized difference maps: mean RGB
difference, grayscale difference, gradient-magnitude difference, and
$1-\mathrm{SSIM}$. Their pixel-wise maximum is
\begin{equation}
M_j(u,v)=\max\left\{D^{\mathrm{rgb}}_j,D^{\mathrm{gray}}_j,
D^{\nabla}_j,D^{\mathrm{ssim}}_j\right\}.
\end{equation}
The primary baseline follows the handout literally and selects one globally
closest template:
\begin{equation}
j^\star=\arg\min_j \operatorname{mean}_{u,v}M_j(u,v),\qquad D=M_{j^\star}.
\end{equation}
The selected map is smoothed with a Gaussian kernel ($\sigma=4$),
morphologically cleaned, and summarized by its 99.5th percentile. The
teammate-inspired per-pixel mode $D(u,v)=\min_j M_j(u,v)$ is retained only as
an ablation; using it as the default would change the meaning of ``closest
template'' and therefore requires explicit disclosure.

\subsection{AnomalyDINO Reproduction}
DINOv2 converts each image into $N$ patch tokens
$\mathbf{z}_i\in\mathbb{R}^{384}$ at the last transformer block. For the
official reproduction, all tokens are L2-normalized and stored in a memory bank
\begin{equation}
\mathcal{M}=\{\bar{\mathbf{z}}_i \mid i\in\mathcal{R}_{c,s,r}\}.
\end{equation}
A PCA mask suppresses background patches based on the first principal component
of the current image. For each retained test patch $\mathbf{x}_q$, we compute
\begin{equation}
d_q=\frac{1}{2}\min_{\mathbf{m}\in\mathcal{M}}
\left\|\bar{\mathbf{x}}_q-\mathbf{m}\right\|_2^2
=1-\cos\!\left(\bar{\mathbf{x}}_q,\mathbf{m}^\star\right).
\end{equation}
The patch distances are reshaped to the token grid. Following AnomalyDINO, the
image score is the mean of the top one percent of patch distances. This
aggregation is robust to isolated noisy patches while retaining sensitivity to
small defects.

The official agnostic preprocessing augments each reference with eight
rotations and applies PCA masking to test images. At $448$ pixels,
ViT-S/14 produces a $32\times32$ patch grid before masking. This coarse grid is
the main source of localization error for thin defects.

\subsection{Multi-Layer Feature Fusion}
Late transformer blocks contain semantic information but discard much
high-frequency spatial detail. Our improvement combines block 9 and block 11:
\begin{equation}
\mathbf{f}_i =
\frac{1}{2}\left(
\frac{\mathbf{z}^{(9)}_i}{\|\mathbf{z}^{(9)}_i\|_2}
+
\frac{\mathbf{z}^{(11)}_i}{\|\mathbf{z}^{(11)}_i\|_2}
\right).
\end{equation}
The fused vectors are again L2-normalized for cosine nearest-neighbor search.
Everything else is unchanged: the same memory size, rotations, PCA mask,
score aggregation, checkpoint, and evaluation code. This isolates the effect of
multi-layer fusion. We expect the earlier layer to add boundary and texture
sensitivity, while the final layer preserves stable object semantics.

\subsection{Avoiding Reference and Threshold Leakage}
For every category, the first 20 lexicographically sorted normal training files
are reserved for threshold calibration. They never enter the memory bank or
few-shot reference pool. Seed $r$ selects contiguous references after this
prefix:
\begin{equation}
\mathcal{R}_{c,s,r} =
\left\{\pi_{20+rs}(c),\ldots,\pi_{20+rs+s-1}(c)\right\},
\end{equation}
where $\pi(c)$ is the sorted training list. Exact filenames are saved for every
run. The demo threshold is the 95th percentile of scores on the calibration
images. AUROC uses unthresholded scores; scores are not rescaled per test image.
